In this chapter, we shall explore the definitions of higher $K$-theory.
The primary goal of higher $K$-theory is to (functorially) define a `$K$-Theory space' $K(R)$ associated to a unital ring $R$.
We then define the higher $K$-groups to be $K_n(R) = \pi_nK(R)$, and verify 
that this description matches with the classical theory for $n = 0, 1, 2$.

A number of approaches can be taken in different contexts, and we have some results showing that in 
some restricted contexts, these various constructions are equivalent, i.e., yield homotopy-equivalent $K$-theory spaces.

Finally, this allows us to define and relate the higher $K$-theory of finitely generated projective modules over a ring, among other applications. 
%Once we have a homotopy-theoretic description of $K$-theory, we can implement tools from Homotopy Theory such as Spectral Sequences, etc.
%towards computations of $K$-theory.

In Section \S1, we first recall some homotopy-theoretic notions.
In Section \S2, we describe the $+$-construction, introduced in [] by <> and used by Quillen in [] to describe higher $K$-groups.
Section \S3 describes higher $K$-groups of a ring using the $+$-construction.
In Section \S4, we describe the $S^{-1}S$ construction on a symmetric monoidal category,
and show its equivalence with the $+$-construction for $(\mathbf{P}(R), \oplus)$.
In Section \S4, we describe Quillen's $Q$-construction for exact categories. 
Section \S5 is devoted to the $+=Q$ theorem, which implies equivalence of the $+$ and $Q$-construction for $\mathbf{P}(R)$.

\section{Homotopy Theory Preliminaries}

A based space is a topological space $X$ along with a distinguished basepoint which we denote $*_X$.
A map of based spaces $X \to Y$ is a map of the underlying topological spaces which preserves basepoints.

We can functorially associate homotopy groups $\pi_n$ to a based space $X$, denoted $\pi_n(X)$. We assume familiarity with the definitions of homotopy and homology groups. 
$H_*(X;M)$ will refer to the singular homology of $X$ with coefficients in some abelian group $M$, or alternatively local coefficients, i.e. $M$ being a $\pi_1(X)$-module.
When $M = \Z$, this is referred to as integral homology. See [] for a detailed treatment.

A CW complex (or cellular complex) $X$ is a colimit of a filtration of $n$-skeleta, denoted $\emptyset = X_{-1} \subseteq X_0 \subseteq X_1 \subseteq X_2 \subseteq \cdots$
where $X_0$ is a discrete set of points and $X_{n}$ is a pushout of $X_{n-1}$ via a disjoint union $\cup_\alpha \varphi_\alpha$ of attaching maps $\varphi_\alpha : S^{n-1}_\alpha \to X_{n-1}$ and
the disjoint union $\cup_\alpha j_\alpha$ of inclusions $j_\alpha : S^{n-1}_\alpha \to D^{n}_\alpha$. (Note $S^{-1}$ is taken to be $\emptyset$.)

That is, $X_{n} = X_{n-1} \cup_{S^{n-1}_\alpha} D^{n}_\alpha$. 
The images of the interiors of the disks $D^n_\alpha$ in $X_n$ are called $n$-cells.
Note that $0$-cells correspond to a set of points, $1$-cells correspond to line segments, $2$-cells to disks, and so on.

The colimit carries the weak topology.
The homology of CW complexes is particularly tractable. See [Hatcher] for a detailed treatment.

\subsection{Geometric Realisation of a small category}
Given a small category $C$, there is a way to interpret $C$ as a CW complex by the following procedure:

\begin{definition}[Simplicial sets]
A simplicial set $X$ is a contravariant functor $\Del \to \Sets$, where $\Del$ is the simplex
category consisting of objects $\mathbf n = \{0,1, \ldots, n-1\}$ and morphisms $\mathbf i \to \mathbf j$ given by weakly-order preserving maps.
Or alternatively, a sequence of sets $X_0, X_1, X_2, \ldots$ with face maps $\partial^n_i : X_n \to X_{n-1}$ and degeneracy maps $\sigma^n_j : X_n \to X_{n+1}$ subject to certain identities $(0 \leq i,j, \leq n)$.
\end{definition}
The category of simplicial sets is denoted $\sSet$. It is clear $\sSet$ is the functor category $[\Del^{\mathord{op}}, \Sets]$.

\begin{example}
The category $\Delta$\footnote{not identical to $\Del$} of standard $n$-simplices $\Delta^n$ with morphisms being the standard inclusion maps $\Delta^i \to \Delta^j$ for $0 \leq i \leq j$
gives a simplicial set by $\mathbf n \to \Delta^n$ and $f : \mathbf i \to \mathbf j$ associated to $\Delta^i \to \Delta^j$.
In fact, this is an isomorphism of categories $\Phi : \Del \cong \Delta$.
\end{example}

\begin{example}
Any space $X$ determines a simplicial set by the representable functor $\Top(\Phi(-),X)$.
This construction is clearly a (covariant) functor $S : \Top \to \sSet$. It is called the \emph{singular functor}.
\end{example}

A simplicial set $X$ may be functorially turned into a topological space, and this functor is called Geometric Realization.
\begin{proposition}[Geometric Realization]
There exists a functor $|-| : \sSet \to \Top$ which is left-adjoint to the singular functor. 
\[
\Top(|X|,Y) \cong \sSet(X, SY)
\]
naturally in $X$ and $Y$.
\end{proposition}
We can define $|-|$ to be the functor taking $X$ to the coend $\int_{\mathbf n \in \Del} X_n \times \Delta^n$ or equivalently,
the colimit $\varinjlim\limits_{\Delta^n \to X} \Delta^n$ in the slice category $\Delta \comma X$.
The result follows due to Yoneda's Lemma. See [Maclane].
Furthermore, $|X|$ has a natural CW structure induced by CW structures on $\Delta^n$. 

\begin{definition}[Nerve of a small category]
Let $C$ be a small category. Then its nerve $NC$ is a simplicial set defined by
taking the $n$-simplices as functors $c : \mathbf{n} \to C$, i.e. diagrams of the form 
\begin{align*}
  c_0 \to c_1 \to \cdots \to c_{n-1} 
\end{align*}
for objects $c_0, c_1, \ldots, c_{n-1}$ in $C$.
The $i$-th face map $\partial^n_i(c)$ on this simplex is obtained by deleting $c_i$ in the evident way.
Similarly, the $j$-th degeneracy map $\sigma^n_j(c)$ is obtained by replacing $c_j$ by $c_j \to c_j$ in the obvious manner.
\end{definition}

We define the \emph{classifying space} of a small category $C$, denoted $BC$, as the geometric realization of $NC$.
\begin{definition}[Classifying space of a small category]
  $BC := |NC|$.
\end{definition}

\begin{example}
For the categories $\mathbf n$, the geometric realization $B \mathbf n = \Phi \mathbf n$.
\end{example}
Furthermore, for any small category $C$, $BC$ is the colimit (with the weak topology) of 
the directed system of the functors $B \mathbf n \to BC$ coming from $\mathbf n \to C$ and the   
morphisms coming from factorizations $\mathbf i \to \mathbf n \to C$.

\begin{proposition}\label{Bfunctorial}
A functor $F : C \to D$ of small categories determines a cellular map $BF : BC \to BD$. Further,
$B(\id_C) : BC \to BC$ is the identity map on $BC$, and for functors $F : C \to D, G : D \to E$, we have
$B(G \circ F) = BG \circ BF$.
\end{proposition}
\begin{proof}
Due to the functoriality of geometric realization on simplicial sets, it suffices to show that
the association $N : \Cat \to \sSet$ is functorial. But this follows directly from axioms of a category.
\end{proof}

Thus we have a functor $B : \Cat \to \Top$.
%\begin{proposition}
%Sucomplexes of $BC$ are in 1-1  with subcategories of $C$.
%\end{proposition}

\begin{example}
A group $G$ may be interpreted as a small category with a single object and all invertible arrows.
We call $BG$ the classifying space of the group $G$.
\end{example}

\begin{proposition}
The fundamental group $\pi_1(BG)$ is naturally isomorphic to $G$. 
%So we have $H_1(BG;\Z) = 0$ if and only if $G$ is perfect.
\end{proposition}

\subsection{$H$-spaces and Homotopy Fibers}
We will note that the following homotopy-theoretic generalization of a topological group will be relevant when considering upcoming constructions.

\begin{definition}[$H$-space]
A $H$-space \footnote{The $H$ stands for Hopf} $X$ is a topological space with a binary operation $\mu : X \times X \to X$ and a distinguished element
$e$ such that 
\begin{enumerate}
  \item $\mu(e,e) = e$ 
  \item $\mu(e,-) \eqv \id_X$ 
  \item $\mu(-,e) \eqv \id_X$ 
\end{enumerate}
\end{definition}
Additionally, we say that a $H$-space is homotopy-associative or homotopy-commutative if those respective properties hold up to homotopy.

Topological groups are trivially examples of $H$-spaces. A class of illuminating nontrivial examples come from $n$-th loop spaces $\Omega^n X$ of some based space $X$, where the multiplication is
given by concatenation of paths.

We refer the reader to [Whitehead] for basic facts about $H$-spaces, some of which we state.

Similar to group completions of monoids, it is a sensible question to ask for a universal 
way to complete a homotopy-commutative homotopy-associative $H$-space into having an \emph{homotopy inverse}. 

Suppose $X$ is a homotopy-commutative homotopy-associative $H$-space. Then its set of connected components form
an abelian monoid, and $H_0(X;\Z)$ is the monoid ring $\Z[\pi_0(X)]$.
Its integral homology $H_*(X;\Z)$ is an associative graded-commutative ring via the Pontryagin product.

\begin{definition}
 A \emph{group-like} $H$-space $X$ is a homotopy-associative homotopy-commutative with
 a homotopy inverse, that is a map $\nu : X \to X$ such that the composites 
 \[
  x \mapsto \mu(x, \nu(x)) \text{ and } x \mapsto \mu(\nu(x), x)
 \]
 are homotopic to the constant map at $e$.
\end{definition}

\begin{theorem}
Let $X$ be a homotopy-associative, homotopy-commutative $H$-space. A \emph{group completion}
of $X$ is a $H$-space map $X \to Y$ such that $\pi_0(X) \to \pi_0(Y)$ is a group completion of the abelian monoid $\pi_0(X)$,
and the natural map from the localization of the homology ring
\[\pi_0(X)^{-1}H_*(X;k) \to H_*(Y;k)\]
is an isomorphism, for any commutative ring $k$.
\end{theorem}


\subsubsection{Homotopy Fiber and Fiber Sequences}

We now describe the important notion of a \emph{homotopy fiber}, which is useful for organizing how homotopy groups behave across a map $f : E \to B$.

\begin{definition}[Homotopy Fiber]
Let $f : E \to B$ be a map of based spaces. The homotopy fiber $F(f)$ of this map is the homotopy-pullback of the cospan $E \to B \leftarrow PB$,
where $PB$ is the (based) space of paths in $Y$.
\end{definition}
More concretely, we can define $F(f)$ to be consisting of pairs $(e,\gamma)$ where $e \in E$ and $\gamma$ is a path in $B$ starting at the basepoint such that 
$f(e) = \gamma(1)$. $F(f)$ inherits the subspace topology from $E \times PB$ where $PB$ itself carries the compact-open topology. We may declare $F(f)$ to have the basepoint $(*_E,*_{PB})$.

The key property of the homotopy fiber is that the sequence $F(f) \to E \to B$ of based space maps which yields a long exact sequence:
\begin{align*}
  \cdots \to \Omega^{n+1} B &\xrightarrow{\partial} \Omega^n F(f) \to \Omega^n E \to \Omega^n B \xrightarrow{\partial} \cdots \\
  \cdots &\xrightarrow{\partial}  \Omega F(f) \to \Omega E \to \Omega B \xrightarrow{\partial} F(f) \to E \to B
\end{align*}

which in turn yields a long exact sequence of groups/pointed sets:
\begin{align*}
  \cdots \to \pi_{n+1} B &\xrightarrow{\partial} \pi_n F(f) \to \pi_n E \to \pi_n B \xrightarrow{\partial} \cdots \\
  \cdots &\xrightarrow{\partial} \pi_1 F(f) \to \pi_1 E \to \pi_1 B \xrightarrow{\partial} \pi_0 F(f) \to \pi_0 E \to \pi_0 B
\end{align*}

Any other such sequence $F \to E \to B$ which results in the evident map $F \to F(f)$ (use $*_{PB}$ in the image) being a homotopy equivalence,
is called a \emph{homotopy fibration sequence}, and consequently also yields such a long exact sequence.

If $f : E \to B$ were a homotopy equivalence, then the homotopy fiber $F(f)$ is weakly contractible, i.e.,
it has the homotopy groups of a point, $\pi_*F(f) \cong \pi_*(\star)$.
If $E$ and $B$ have homotopy type of CW complexes, then so is $F(f)$ and by Whitehead's Theorem [], weak homotopy equivalence implies homotopy equivalence.

Conversely, if the fiber $F(f)$ is contractible, then $f$ is a weak homotopy equivalence, i.e., the induced maps $\pi_n(E) \to \pi_n(B)$ are isomorphisms.
If $E$ and $B$ have homotopy type of CW complexes, then $f$ is a homotopy equivalence.

\subsection{Homotopy Theoretic Aspects of Categories}

Owing to the geometric realization functor $B : \Cat \to \Top$, we can 
think of small categories as topological spaces via $B$.
This allows us to extend topological notions to categories by viewing them through their geometric realization. 

By abuse of notation, we will make statements such a category $C$ is contracible, or connected, etc. (any homotopy-theoretic property)
if its geometric realization is so. 
If a functor $F : C \to D$ induces a homotopy equivalence $BF : BC \to CD$, then we say that $F : C \to D$ is a homotopy equivalence,
and the categories $C$ and $D$ are homotopy equivalent.

\begin{proposition}
A natural transformation $\eta : F \Rightarrow G$ between functors $F, G : C \to D$
describes a homotopy $BC \times [0,1] \to BD$ which we denote $B\eta : BF \To BG$.
\end{proposition}
\begin{proof}
We can view $\eta : F \To G$ as a functor $C \times \mathbf 2 \to D$ whose restriction
to $C \times {0}$ is $F$ and the restriction to $C \times {1}$ is $G$. 
Upon taking the geometric realization, we are done.
\end{proof}

A consequence is that an adjoint pairs of functors $L : C \to D, R : D \to C$ 
with natural transformations $LR \To \id_D$ and $\id_C \To RL$, leads to a 
homotopy equivalence $BC \eqv BD$. In particular, 

\begin{corollary}
An equivalence of small categories is a homotopy equivalence.
\end{corollary}

\begin{remark}
The converse is not true. If $F : C \to D$ is a homotopy equivalence, then it need not be an equivalence of categories
in the usual sense. Take for instance  [][]
\end{remark}

\begin{proposition}
If a small category $C$ has an initial or a final object, then $BC$ is contractible.
\end{proposition}
\begin{proof}
Let $C$ have an initial object $c_0$, then the (unique) functor $F : C \to \mathbf 1$
has a left adjoint  $L : \mathbf 1 \to C$ given by $0 \mapsto c_0$.
Dually, if $C$ has a final object $c_1$, then $F$ has a right adjoint $R : \mathbf 1 \to C$ given by $0 \mapsto c_1$.
The result follows from the remark made earlier.
\end{proof}

\subsubsection{Fundamental Group of a Category}

\subsection{Quillen's Theorems A and B}

In this section, we prove two theorems originally proven by Quillen in [Quillen 1973],
which shall be integral towards proving the $+=Q$ Theorem in [].

Theorem A gives sufficient conditions for a functor $F : C \to D$ of small categories to be 
a homotopy equivalence.
Theorem B gives sufficient conditions to create homotopy fibration sequences from
functors $F : C \to D$.

But first, we introduce the notion of a bisimplicial set and some of its properties.

\begin{definition}
A \emph{bisimplicial set} $X$ is a contravariant functor $\Del \times \Del \to \Sets$.
Alternatively, it is a doubly indexed family $X_{p,q}$ of sets, together with
\emph{horizontal} face and degeneracy maps
\[\partial^h_i : X_{p,q} \to X_{p-1,q} \text{ and } \sigma^h_i : X_{p,q} \to X_{p+1,q}\]
and \emph{vertical} face and degeneracy maps
\[\partial^v_i : X_{p,q} \to X_{p,q-1} \text{ and } \sigma^v_i : X_{p,q} \to X_{p,q+1}\]
such that the horizontal maps commute with the vertical maps, and both pairs follow
the required relations among face and degeneracy maps as for simplicial sets.
In particular, each $X_{\bullet,q}$ and $X_{p,\bullet}$ defines a simplicial set.
See [] for a detailed treatment of bisimplicial sets.
\end{definition}

\begin{definition}
The \emph{geometric realization} of a bisimplicial set $X$
is the cellular complex $BX$ given by taking the disjoint union of $X_{p,q} \times \Delta^p \times \Delta^q$
and identifying its horizontal and vertical faces with the appropriate
$\Delta^{p-1} \times \Delta^q$ or $\Delta^p \times \Delta^{q-1}$.

Alternatively, it is the coend $\int_{p,q \in \Del} X_{p,q} \times \Delta^p \times \Delta^q$.
\end{definition}

Note there is the \emph{diagonal} functor from bisimplicial sets to simplicial sets
$\diag(X) \mapsto X_{n,n}$, and it is well-known that $BX$ is homeomorphic to $B \diag(X)$.
See for instance, [Bousfield]. The following theorem is also well known: (see [Wald 78])

\begin{theorem}\label{thm:bisimplicial}
 Let $f : X \to Y$ be a map of simplicial sets. 
 \begin{enumerate}
    \item If each simplicial map $X_{p,\bullet} \to Y_{p,\bullet}$ is a homotopy equivalence,
    so is $BX \to BY$.
    \item If $N$ is the nerve of a category $C$, though of as a bisimplicial set which is constant
    in the second coordinate,
    and $f^{-1}(c, _\bullet) \to f^{-1}(c', _\bullet)$ is a homotopy equivalence for every $c \to c'$ in $C$,
    then each $B(f^{-1}(c, _\bullet)) \to BX \to BC$ is a homotopy equivalence.
 \end{enumerate}
\end{theorem}

Using this, we prove a lemma which eventually helps us prove Theorems A and B.
\begin{lemma}[Quillen] \label{lemma:bisimplicial}
If $F : C \to D$ is a functor, the evident functor $D \comma F \to C$ given by mapping $(d, d \to Fc)$ to $c$
is a homotopy equivalence.
\end{lemma}
\begin{proof}
Let $X$ denote the bisimplicial set where $X_{p,q}$ is the set of all pairs of sequences
\[
(d_q \to \cdots \to d_0 \to F(c_0), c_0 \to \cdots \to c_p)
\]
with the horizontal and vertical faces inherited from the nerves of $C$ and $D$.
Consider the projection $X \to NC$ which preserves only the second sequence in the pair.
Note that the discrete set $NC_p$ is the set of all sequecnces $c_0 \to \cdots \to c_p$,
and its inverse image is isomorphic to the nerve of $D \comma F(c_0)$.
$D \comma F(c_0)$ is contractible since it has a terminal object $(c_0,\id_{F(c_0)})$.
Using Theorem \ref{thm:bisimplicial}, from the homotopy fibration sequence we get $BX \xlongrightarrow{\eqv} BC$.
The simplicial set $\diag(X)$ will be isomorphic to the nerve of $D \comma F$, and thus we have the composite
$B(D \comma F) \xlongrightarrow{\eqv} BX \xlongrightarrow{\eqv} BC$, which concludes the proof.

\end{proof}

Now, we are in a position to prove Theorem A. 

\begin{theorem}[Quillen's Theorem A]\label{thm:A}
Let $F : C \to D$ be a functor of small categories. If the category $d \comma F$ is contractible
for every object $d$ in $D$, then $F$ is a homotopy equivalence. 
\end{theorem}
\begin{proof}
Consider the comma category $D \comma F$ and the evident functors $D \comma F \to C$
and $D \comma F \to D^{op}$ given by mapping $(d, d \to Fc)$ to $c$ and $d$ respectively.
Note that $B(D\comma F) \to BC$ is a homotopy equivalence by \ref{lemma:bisimplicial}.

We have a functor $D\comma F \to D\comma\id_D$ given by mapping $(d, d \to Fc)$ to $(d, \id_d)$,
and a functor $D \comma \id_D \to D^{op}$ given by mapping $(d,d \to Fc)$ to $d$. The latter is a homotopy equivalence
since it is an equivalence of categories.

But we have a commutative diagram of categories 
\[
\begin{tikzcd}
D^{op} \arrow[d, equal] & D\comma F \arrow[l] \arrow[d] \arrow[r] & C \arrow[d] \\
D^{op} & D \comma \id_D \arrow[l] \arrow[r] & D
\end{tikzcd}
\]
So it suffices to show that $D \comma F \to D^{op}$ is a homotopy equivalence.
Let $X$ be the bisimplicial set of Lemma \ref{lemma:bisimplicial}. Then the map $B(D \comma F) \xrightarrow{\pi} B D^{op}$
factors as $B(D \comma F) \eqv BX \to BD^{op}$. The image of the simplicial map $\pi_{\bullet,q} : X_{\bullet,q} \to (ND^{op})_q$
is the discrete set of all sequences $d_q \to \cdots \to d_0$ in $D$.
For any such sequence, the inverse image in $X_{\bullet, q}$ is isomorphic to the nerve of $d_0 \comma F$.
Since we assumed that the latter to be contractible, by Theorem \ref{thm:bisimplicial},
we have the desired homotopy equivalence $B(D \comma F) \xlongrightarrow\eqv BX \xlongrightarrow\eqv BD$.
\end{proof}

And then we have Theorem B, which will later be used to construct homotopy fibration sequences of categories.

\begin{theorem}[Quillen's Theorem B]\label{thm:B}
 Let $F : C \to D$ be a functor of small categories such that every morphism $d \to d'$ in $D$
 induces a homotopy equivalence $d \comma F \to d' \comma F$. Then for any object $d$ of $D$,
 the geometric realization of the sequence
 \[
 d \comma F \to C \to D
 \]
 is a homotopy fibration sequence. Consequently, we have a long exact sequence 
 \[
 \cdots \to \pi_{n+1}(BD) \xrightarrow{\partial} \pi_n B(d \comma F) \xrightarrow{j} \pi_n BC \xrightarrow{F}\pi_n BD \to \cdots
 \]
\end{theorem}
\begin{proof}
Consider the projection $\pi : X \to N(D^{op})$ as in Theorem A (\ref{thm:A}).
Since the inverse image $\pi^{-1}(d)$ for any $d \in D$ is the nerve of $d \comma F$
by the same argument, and we have the condition that $d \to d'$ induces $d \comma F \to d' \comma F$ is 
a homotopy euqivalence, we can use Theorem \ref{thm:bisimplicial} to deduce that
$B(d \comma F) \to BX \to BD^{op}$ is a homotopy fibration sequence.

Since $BX \cong B(\diag X) \cong B(D \comma F)$, we have $B(d \comma F) \to B(D \comma F) \to BD^{op}$ 
is a fibration sequence.
Similarly, $d \comma \id_D \to D \comma \id_D \to D^{op}$ is also a fibration sequence.
Consider the diagram
\[
\begin{tikzcd}
D^{op} \arrow[d, equal] & D\comma F \arrow[l] \arrow[d, "F"] & d \comma F \arrow[d] \arrow[l] \\
D^{op} & D \comma \id_D \arrow[l, "\eqv"'] & d \comma \id_D \arrow[l] 
\end{tikzcd}
\]
Note the $d \comma \id_D$ is contractible. So $D \comma \id_D \xrightarrow{\eqv} D^{op}$ is an equivalence,
giving us the fibration sequence $B(d \comma F) \to B(D \comma F) \to B(D^{op}) = BD$.
The theorem follows from the fact that 
$j : d \comma F \to BC$ induces $B(D \comma F) \xrightarrow{\eqv} BC$.

\end{proof}
\subsection{$+$-construction on a connected CW complex}

The $+$-construction, devised by <> in [], is a universal way to kill a perfect normal subgroup of the fundamental group of a connected 
CW complex, while preserving homology theory. 
The discussion of the $+$-construction fits into the framework of acyclic maps, which we explore here.

\begin{definition}
A space $X$ is said to be acyclic if it has the homology of a point, i.e., $H_*(X;\Z) \cong H_*(\star;\Z) = \underline{\Z}$.
\end{definition}
In particular, contractible spaces are acyclic. Acyclic spaces serve an analogous purpose to homology as contractible spaces do to homotopy.

\begin{definition}[Acyclic maps]
Let $f : X \to Y$ be a map of based connected CW complexes. We say $f$ is acyclic if the homotopy fiber $F(f)$ is acyclic, i.e., has the homology of a point. 
\end{definition}
In particular, a homotopy equivalence is acyclic since the former implies that the homotopy fiber is contractible.
Note an acyclic map induces isomorphisms on integral homology $f_* : H_*(X;\Z) \cong H_*(Y;\Z)$.
Furthermore, we have the following result characterising acyclic maps:

\begin{lemma}\label{lemma:acyclic}
Let $X, Y$ be based connected CW complexes. A map $f : X \to Y$ is acyclic if and only if for every $\pi_1(Y)$-module
$M$, it induces an isomorphism $H_*(X;M) \cong H_*(Y;M)$.
\end{lemma}
The proof for this uses Serre Spectral sequences, which we skip and refer the reader to [Weibel].

\begin{definition}[relative $+$-construction]
Let $X$ be a based connected CW complex and $P$ be a perfect normal subgroup of $\pi_1(X)$. 
Then we say a based acyclic map of $CW$-complexes $f : X \to Y$ is a $+$-construction
relative to $P$ if $P$ is the kernel of the induced map $\pi_1(X) \to \pi_1(Y)$.
\end{definition}

Since $f$ was assumed to be acyclic, we can deduce from the exact sequence $\pi_1 F(f) \to \pi_1(X) \to \pi_1(Y) \to \pi_0 F(f)$ that
the map $\pi_1(X) \to \pi_1(Y)$ is onto. (because $\pi_0 F(f) = \{\star\}$)
And $\pi_1(Y) \cong \pi_1(X)/P$.

%Note that since $F(f)$ is acyclic, the exact sequence $\pi_1 F(f) \to \pi_1 X \to \pi_1 Y \to \pi_0 F(f)$ gives us that 
%$\pi_1(X) \to \pi_1(Y)$ is surjective
\begin{theorem}[Quillen]\label{thm:plusexists}
Let $X$ be a based connected CW complex and $P$ be a perfect normal subgroup of $\pi_1(X)$.
\begin{enumerate}
  \item There exists a $+$-construction $X \to Y$ relative to $P$.
  \item A $+$-construction $f : X \to Y$ is universal among maps which kill $P$ in the fundamental group.
  That is, for any $g : X \to Z$ such that the image of $P$ in $\pi_1(Z)$ vanishes, we have a map $h : Y \to Z$, unique up to homotopy, such that $g = hf$. 
  \item In particular, if $g$ is another $+$-construction relative to $P$, then there is a homotopy equivalence $h : Y \xlongrightarrow{\sim} Z$ compatible with the maps.
\end{enumerate}
\end{theorem}
The proof of this uses topological obstruction theory. We describe a construction for $1$ in \S\ref{ss:construction}
and refer the reader to [Berrick] for a full proof of 1 and 2. 3 follows immediately from 2 since we have maps both way which must compose to identity up to homotopy.

\begin{remark}
Any group $G$ has a maximal perfect normal subgroup, which we call the \emph{perfect radical}. In fact, it is the subgroup generated by all perfect normal subgroups,
hence their maximum.
From now on, for any based CW complex $X$, we will denote by $X^+$ any $+$-construction on $X$ relative to the perfect radical $P \triangleleft \pi_1(X)$.
\end{remark}

\subsection{Explicit Construction} \label{ss:construction}
A number of explicit \emph{models} exist for the $+$-construction. Here, we present a construction due to [Kervaire] which constructs $X^+$ by attaching $2$-cells 
to kill the maximal perfect subgroup of $\pi_1(X)$, and then $3$-cells to recover the integral homology of $X$, making sure
that the inclusion $X \hookrightarrow X^+$ is acyclic.
We refer the reader to [Gersten] and [Suslin] for a variety of constructions. 
We will describe more models for $K$-theory in Sections \S3 and \S4.

\subsubsection*{Construction}
Let $X$ be a connected CW complex with a perfect normal subgroup $P$ of the fundamental group $\pi_1(X)$.
Identify the $1$-cells in $X$ corresponding to of $P$ as $e^1_p$ for each $p \in P$. 
Now construct a pushout $Y$ by attaching disks $D^2_p$ with the boundary $S^2_p$ along every $e^1_p$.
Call these new cells $e^2_p$.

That is, $X'$ is the pushout of the diagram 
\[
\begin{tikzcd}%[column sep=large]
  \coprod\limits_{p \in P} S^1_{p} \arrow[r, hook] \arrow[d] & \coprod\limits_{p \in P} D^2_{p} \arrow[d] \\
  X \arrow[r, hook] & Y
\end{tikzcd}
\]

Using the van-Kampen Theorem, we have the pushout through $\pi_1$ in $\mathbf{Grp}$,
which gives us
  $\pi_1(Y) = \pi_1(X) *_{P} 1
         = \pi_1(X)/P$.
         (amalgamated product)

Consequently, $H_1$ is the abelianization. Since $P$ is perfect, it is killed by abelianization and we have
$ H_1(Y) = H_1(X)$.

Consider the pair $(Y,X)$. Since only $2$-cells were attached, the homology of the pair is
\[
H_2(Y,X) = \bigoplus_{[e^2_p]} \Z
\]
and $0$ elsewhere.
We wish to cleverly attach $3$-cells to kill relative homology without affecting $H_3$.

Consider the long exact sequence of the pair $(Y,X)$
\begin{align*}
  \cdots \to H_{n+1}(Y,X) \xrightarrow{\partial} H_{n}(X) &\to H_n(Y) \to H_n(Y,X) \xrightarrow{\partial} \cdots 
\end{align*}

At $n=1$, since the map $H_1(X) \to H_1(Y)$ is an isomorphism, we have that $H_2(Y,X) \rightarrow{} H_1(X)$ is $0$.
Consequently, the map $H_2(Y) \to H_2(Y,X)$ is onto.
We can lift each representative in $H_2(Y,X)$ to a representative in $H_2(Y)$, say $\sigma_p$ for each class $[e^2_p]$.

Note each $\sigma_p$ is represented by a map $h_p : S^2_p \to Y$.
Now attach $3$-cells $e^3_p$ corresponding to each $\sigma_p$ along the maps $h_p : S^2_p \to Y$.

This yields another pushout which we call $Z$
\[
  Z = Y \cup \bigcup\limits_{[e^2_p]} e^3_p 
  \]

The triple $X \subseteq Y \subseteq Z$ induces a long exact sequence in relative homology
\begin{align*}
  \cdots \to H_{n+1}(Z,Y) \xrightarrow{\partial} H_{n}(Y,X) \to H_{n}(Z,X) \to H_n(Z,Y) \xrightarrow{\partial}\cdots 
\end{align*}

At $n=1$, since $H_2(Y,X) = 0$ and $H_0(Z,Y) = 0$, this gives us
\begin{align*}
  \cdots \xrightarrow{\partial} 0 \to H_{1}(Z,X) \xrightarrow{\cong} H_1(Z,Y)\xrightarrow{\partial} 0 \to  \cdots 
\end{align*}

Similarly, for $n \geq 4$, since $H_{n+1}(Z,Y) = 0$ and $H_{n-1}(Y,X) = 0$, we get
\begin{align*}
  \cdots \xrightarrow{\partial} 0 \to H_{n}(Z,X)\xrightarrow{\cong}H_n(Z,Y) \xrightarrow{\partial} 0 \to \cdots 
\end{align*}

So we conclude $H_n(Z,X) \cong H_n(Z,Y) = 0$ for $n \neq 1,2,3$.

The only nontrivial terms in the sequence that remain are 
\begin{align*}
  \cdots \xrightarrow{\partial} 0 \to H_3(Z,X) \to H_3(Z,Y) \xrightarrow{\partial} H_2(Y,X) \to  0 \to \cdots 
\end{align*}

But $H_3(Z,Y)$ and $H_2(Y,X)$ are computable precisely due to construction.
$Z$ was constructed so as to kill  
\[
H_3(Z,Y) = \bigoplus_{[\sigma_p]} \Z
\]
which is isomorphic via the evident maps to $H_2(Y,X) = \bigoplus_{[e_p]} \Z$.

So finally, we have $H_3(Z,X) = 0$, and altogether $H_n(Z,X) = 0$ for all $n$. 
Consequently, in the long exact sequence of the pair $(Z,X)$, we get
\begin{align*}
  \cdots \to H_{n+1}(Z,X) \xrightarrow{\partial} H_n(X) \to H_n(Z) \to H_n(Z,X) \xrightarrow{\partial} \cdots 
\end{align*}
which gives us that $H_n(X) \cong H_n(Z)$ induced via the inclusions are isomorphisms for all $n$.

It is clear that $P$ is the kernel of the induced map $\pi_1(X) \to \pi_1(Y)$.
This concludes the result that the stated construction preserves integral homology. 

%To extend this to any $\pi_1(Z)$ module $M$, we need to lift the problem to their universal covers.
%Let $\widetilde X \to X$ be the universal cover of $X$.

If this result were to extend to local coefficient systems, this would mean that the map $X \to Z$ is acyclic 
and forms a model for $X^+$. We do not pursue the full proof here and instead refer the reader to [Rosenberg].
Indeed, there is an alternate result in Proposition \ref{prop:acyclic_H_space} which shows that what we have proven suffices
for the case we are interested in, $BGL(R)^+$, which is what we shall discuss now.

\section{The $BGL(R)^+$ construction for a ring $R$} \ref{sec:BGL(R)^+}
We apply the $+$-construction on the classifying space $BGL(R)$ of the stable general linear group over $R$.
The maximal perfect subgroup of $GL(R) = \pi_1BGL(R)$ is the stable group of elementary matrices, $E(R)$.

\begin{theorem}
With this definition of $BGL(R)^+$, we obtain:
\begin{enumerate}
\item $\pi_1BGL(R)^+ \cong K_1(R)$ with the kernel of the induced map $GL (R) = \pi_1BGL(R) \to \pi_1BGL(R)^+$ being $E(R)$.
\item We have induced isomoprhisms $H_*(BGL(R);M) \cong H_*(BGL(R)^+;M)$ for every $K_1(R)$-module $M$.
\end{enumerate}
\end{theorem}
\begin{proof}
Since $E(R)$ is the maximal perfect normal subgroup of $GL(R)$, 1 follows immediately from Theorem \ref{thm:plusexists}. 2 follows from the characterization of acyclic maps in Lemma \ref{lemma:acyclic}. 
\end{proof}

Due to Theorem \ref{thm:plusexists}, $BGL^+(R)$ is uniquely defined up to homotopy equivalence via a canonical equivalence.
Consequently, we may define $K_n(R) = \pi_n BGL(R)^+$ for $n \geq 1$. 
It is obvious that $K_1(R)$ defined here matches the description in Chapter 3.
It remains to verify that this agrees with the classical theory for $n = 2$.

\begin{proposition}
Let $P$ be a perfect normal subgroup of $G$ with a $+$-construction $f : BG \to BG^+$ relative to $P$.
Let $F(f)$ be the homotopy fiber of $f$. Then $\pi_1F(f)$ is the universal central extension of $P$, and consequently $\pi_2 (BG^+) \cong H_2(P;\Z)$. 
\end{proposition}
\begin{proof}
We have the fibration long exact sequence
\[
  \cdots \to \pi_2(BG) \to \pi_2(BG^+) \to \pi_1 F(f) \to G \to G/P \to 1
\]
Since $BG$ is an Eilenberg-Maclane space $K(G,1)$, we have $\pi_2(BG) = 0$. 
So $\pi_2(BG^+)$ injects into $\pi_1F(f)$, which in turns surjects onto the kernel of $G \to G/P$ which is $P$.
Since $\pi_2(BG^+)$ is abelian, it must inject into the center of $\pi_1 F(f)$.
This tells us that $\pi_1F(f) \to P$ is a central extension.
Since $f$ is acylic, we have $\pi_1 F(f)$ is perfect and so $H_2(\pi_1 F(f), \Z) = 0$.
By the Recognition Theorem [Chapter 3], we have that $\pi_1F(f)$ is the universal central extension of $P$.
Consequently, the kernel $\pi_2(BG^+)$ is given by $H_2(P;\Z)$.

\end{proof}

\begin{corollary}
We have $\pi_2(BGL(R))^+ \cong H_2(E(R);\Z)$, and thus our definition matches with the classical theory
since $H_2(E(R);\Z) \cong K_2(R)$ from Chapter 3. 
\end{corollary}

However, at this point, it is not clear that the associations $R \mapsto BGL(R)^+$ or $R \mapsto K_n(R)$ are functorial.
This invites some discussion.

\subsection{Functoriality of the $BGL(R)^+$-construction}
Apart from the upcoming models via $S^{-1}S$ and $Q$-constructions, we can make the following argument to show that
the constructions of $BGL(R)^+$ can be performed functorially.

We can first define $BGL(\Z)^+$ and take a homotopy pushout (using the double mapping cylinder) via the canonical map $\Z \to R$ which induces $BGL(\Z) \to BGL(R)$.
Note that the homotopy pushout 
\[
\begin{tikzcd}
  BGL(\Z) \arrow[r] \arrow[d] & BGL(\Z)^+ \arrow[d] \\
  BGL(R) \arrow[r] & BGL(R)^+
\end{tikzcd}
\]
is a plus construction for $BGL(R)$ and is functorial in $\Z \to R$.

\subsection{$H$-Space Properties of $BGL(R)^+$}

It's not hard to see that we can put a $H$-space structure on $BGL(R)^+$. 
\begin{proposition}
$BGL(R)^+$ is a homotopy-associative $H$-space.
\end{proposition}
\begin{proof}
We have the group map $\oplus : GL_m(R) \times GL_n(R) \to GL_{m+n}(R) \subset GL(R)$ 
which induces $BGL_m(R) \times BGL_n(R) \to BGL(R)$. 
Since we have show that the plus construction can be taken to be functorial,
we have maps \[(BGL_m(R) \times BGL_n(R))^+ \to BGL(R)^+\].
[Finish this]
\end{proof}

There is a recognition criteria for $+$-constructions of $BGL(R)$ in terms of homology and $H$-spaces, due to Quillen.
The proof again uses obstruction theory. See [Gersten].

\begin{theorem}
The plus construction $i : BGL(R) \to BGL(R)^+$ is universal for maps into $H$-spaces. That is, for any map
$f : BGL(R) \to H$ for some $H$-space $H$, we have a map $g : BGL(R)^+ \to H$ so that $f = g i$,
such that $g$ is unique up to weak homotopy, i.e. the induced maps $\pi_n(BGL(R)^+) \to \pi_n(H)$ are independent of $g$.
\end{theorem}

\begin{proposition}\label{prop:acyclic_H_space}
If further $f_* : H_*(BGL(R);\Z) \xlongrightarrow{\cong} H_*(H;\Z)$ is an isomorphism, then $f$ is acyclic
and $g$ is a homotopy equivance.
\end{proposition}
%\begin{proof}
%It suffices to prove that $f$ is acyclic.
%\end{proof}

\section{$S^{-1}S$ construction}
For a symmetric monoidal category $(S, \square)$, the construction of Chapter 2 for $K_0$ may be extended to 
a $K$-theory space $K^\square(S)$ by constructing an appropriate group-completion of $S$ with respect to the monoidal operation $\square$.

Note that the geometric realization of a symmetric monoidal category is a homotopy-associative homotopy-commutative $H$-space
with the multiplication coming from the monoidal product $\square$.
We wish to construct a symmetric monoidal category $S^{-1}S$ which describes a group completion of $BS$ so as to make $BS^{-1}S$ be a group-like $H$-space.
[Grayson]

But first, for either of $BS$ or $BS^{-1}S$ to be interesting, we must reduce our scope.

\begin{definition}[symmetric monoidal groupoid]
A \emph{symmetric monoidal groupoid} is a symmetric monoidal category $(S,\square)$ in which all morphisms are isomorphisms.

Further, we say $S$ is \emph{faithful} if all translations $s \square : \Aut(t) \to \Aut(s \square t)$ are injective. 
\end{definition}

\begin{definition}
Given a symmetric monoidal category $S$, $\iso(S)$ is the symmetric monoidal groupoid which is a subcategory of $S$ consisting of the same objects but inheriting only isomorphisms.
\end{definition}
A variety of symmetric monoidal groupoids yield interesting geometric realizations.
\begin{example}
The category $\mathbf F(R)$ with $\square = \oplus$ is contractible since it has an initial object $0$. However $\iso(\mathbf F(R))$ 
is more interesting and in fact equivalent to the disjoint union $\coprod_n GL_n(R)$ where $n$ runs through the integers.
\end{example}
\begin{example}
Similarly, $0$ is also an initial object for the category $\mathbf P(R)$ with $\square = \oplus$.
But $\iso(\mathbf P(R))$ is equivalent to $\coprod_P \Aut(P)$ where $P$ runs over representatives of isomorphism clases 
of finitely generated projective modules.
\end{example}

Let $S$ be a faithful symmetric monoidal groupoid. 

\begin{definition}[$S^{-1}S$]
Define $S^{-1}S$ to be a category where the objects are pairs of objects $(m,n)$ for $m, n \in S$. 
Morphisms in $S^{-1}S$ from $(m_1,m_2) \to (n_1,n_2)$ are equivalence classes of triples $(s,f_1,f_2)$ for some object $s \in S$ and 
$(f_1,f_2) : (s \square m_1, s \square m_2) \to (n_1,n_2)$ being a morphism in $S \times S$.

That is,
\[
\begin{tikzcd}
(m_1,m_2) \arrow[r, "s\square"] & (s\square m_1, s \square m_2) \arrow[r, "(f_1{,}f_2)"] & (n_1, n_2)
\end{tikzcd}
\]
which we consider equivalent to 
\[
\begin{tikzcd}
(m_1,m_2) \arrow[r, "t\square"] & (t\square m_1, t \square m_2) \arrow[r, "(g_1{,}g_2)"] & (n_1, n_2)
\end{tikzcd}
\]
if there is an isomorphism $\alpha : s \xlongrightarrow{\cong} t$ such that $\alpha \square g_1 = f_1$ and $\alpha \square g_2 = f_2$. 
The composition $[(s,f_1,f_2)]\circ[(t,g_1,g_2)]$ is defined as $[(s \square t, f_1 \circ g_1, f_2 \circ g_2)]$.
One may verify this is well-defined and indeed forms a category.
\end{definition}

\begin{remark}
Consider two distinguished classes of morphisms in $S^{-1}S$.
One is simultaneous translations, which we formally write $s \square : (m_1,m_2) \to (s\square m_1, s \square m_2)$, represented by $(s,\id,\id)$. 
The other being pairs of maps $(f_1, f_2) : (m_1, m_2) \to (n_1, n_2)$ coming from morphisms in $S \times S$ where the translation composite is trivial.
That is, the morphism is represented by $(e, f_1, f_2)$.
\end{remark}

\begin{proposition}
Let $S$ be a faithful symmetric monoidal groupoid. Then any morphism $[(s,f_1,f_2)]$ in $S^{-1}S$ 
determines $s, f_1, f_2$ uniquely up to unique isomorphisms.
\end{proposition}
\begin{proof}
It is clear $s, f_1, f_2$ are determined uniquely up to isomorphism.
Suppose $(s,f_1,f_2)$ is equivalent to $(t,g_1,g_2)$ via isomorphisms $\alpha, \beta : s \xlongrightarrow{\cong} t$ 
Since $\alpha \square g_i = f_i = \beta \square g_i$, 
we have $\alpha \square m_i = \beta \square m_i$ and since $S$ was faithful, we recover $\alpha = \beta$.
\end{proof}

Thus, we can think of morphisms $[(s,f_1,f_2)]$ in $S^{-1}S$ as comprising of a \emph{translation component} $[s \square]$ and
\emph{mapping component} $[(f_1,f_2)]$.

\begin{remark}
$S^{-1}S$ is a symmetric monoidal category with the multiplication given by $(m_1,m_2) \square (n_1, n_2) := (m_1 \square n_1, m_2 \square n_2)$.
It is clear that the functor $S \to S^{-1}S$, which sends $s \mapsto (s,e)$ is a monoidal functor.
Further, a monoidal functor $S \to T$ describes a (monoidal) natural transformation
\[
\begin{tikzcd}
S \arrow[r] \arrow[d] & T \arrow[d] \\
S^{-1}S \arrow[r] & T^{-1}T
\end{tikzcd}
\]

We also have an induced $H$-space map $BS \to B(S^{-1}S)$, and a map $\pi_0(S) \to \pi_0(S^{-1}S)$
of abelian monoids. (The component of $s$ maps to the component of $(s,e)$) 

Further, $\pi_0(S^{-1}S)$ is an abelian group since for the connected coponent of $(m,n)$ we have we have inverse $(n,m)$ 
due to the isomorphism $(e,e) \xlongrightarrow{\cong} (m,n) \square (n,m)$
given by the mapping component $[(c_{mn},c^{-1}_{mn})]$ where $c_{mn}$ is the symmetric braiding isomorphism $m \square n \cong n \square m$ in $S$.
It is abelian due to $S$ being symmetric. 

%(Warning: these isomorphisms $(e,e) \to (m,n) \square (n,m)$ are not natural in $m, n$ due to both an isomorphism and its inverse being used in the definition.)
\end{remark}

\begin{definition}
Let $S$ be a symmetric monoidal groupoid. We define 
\[
K^\square(S) := B(S^{-1}S)
\]
and consequently we can define 
$K^\square_n(S) = \pi_n B(S^{-1}S)$.
  
\end{definition}
Note that a (strict) monoidal functor $S \to T$ induces $S^{-1}S \to T^{-1}T$ which in turn induces $K^\square(S) \to K^\square(T)$, 
Thus the $S^{-1}S$ construction is functorial!

\begin{proposition}
The group $K_0^\square(S) = \pi_0(S^{-1}S)$ is indeed the group completion of the monoid given by 
isomorphism classes of $S$ via the inclusion $i_*: \pi_0(S) \to \pi_0(S^{-1}S)$.
\end{proposition}
\begin{proof}
Let $\varphi : \pi_0(S) \to G$ be a group completion of $\pi_0(S)$.
We will show that the universal homomorphism $\widetilde i_* : G \to \pi_0(S^{-1}S)$ is invertible.
Let $j$ be the set map $S^{-1}S \to  G$ as $j(m,n) := \varphi(m) - \varphi(n)$.
If we have $s \in S$ then 
$j(m_1,m_2) = j(s\square m_1, s\square m_2)$ and if $f_i : m_i \to n_i$, then $j(m_1,m_2) = \varphi(m_1) - \varphi(m_2) = \varphi(n_1) - \varphi(n_2) = j(n_1,n_2)$.
This shows $j$ induces a well-defined homomorphism $j_* : \pi_0(S^{-1}S) \to G$. 
By construction, it is the inverse to $G \to \pi_0(S^{-1}S)$.
%We'll show $i_*$ satisfies the universal property. Let $G$ be a group and $\varphi : \pi_0(S) \to G$ be a monoid homomorphism.
%Then define $\widetilde \varphi : \pi_0({S^{-1}S}) \to G$ by mapping the component of $(m,n)$ to $\varphi([m]) - \varphi([n])$.
%Clearly, we have $\widetilde \varphi \circ i_* = \varphi$ since $i_*([m]) = [(m,e)]$.
%Furthermore, this choice is fixed 
\end{proof}
This shows our definition matches the definition of $K_0(S)$ in Chapter 2.

\subsection{Actions on other Categories}
\begin{definition}
A monoidal category $S$ is said to \emph{act upon} a category $X$ by a functor $\boxtimes : S \times X \to X$,
if there are natural isomorphisms 
  \[s \square  (t\boxtimes  x) \cong (s\square t) \boxtimes  x \qquad \text{ and } \qquad e\boxtimes  x \cong x\]
for $s,t \in S$ and $x \in X$,
satisfying the coherence conditions for associativity and unit law, i.e. pentagonal and triangular diagrams of \ref{symmetric_monoidal_category}.
\end{definition}

For instance, $S$ acts on itself by $\square$. If $X$ is a discrete category (a set), then $S$ acts on $X$ exactly when the monoid of isomorphism classes of $S$ 
acts on $X$. 

\begin{definition}[Translation category]
If $S$ acts upon $X$, then the \emph{translation category} $\langle S, X \rangle$ is the category with the same objects as $X$, and morphisms $x \to y$
are equivalence classes of pairs $(s, s \boxtimes x \to y)$ for some $s \in S$.

Two pairs $(s,s \boxtimes x \to y)$ 
and $(t, t \boxtimes x \to y)$ are considered equivalent if there is an isomorphism $s \cong t$ such that the 
composite morphism $t \boxtimes x \cong s \to y$ is the same as $t \boxtimes x \to y$.
\end{definition}

\begin{example}
Given a category $X$ such that $S$ acts on $X$ by $\boxtimes$, we will write $S^{-1}X$ for the category $\lr {S, S \times X}$ with the action
given by $s \boxtimes (t,x) = (s \square t, s \boxtimes x)$.
Note that if $X = S$, then this recovers $S^{-1}S$.
\end{example}

\begin{remark}\label{translation contractible}
If $S$ were a monoidal groupoid (every arrow in $S$ is an isomorphism), then $e$ would be an inital object of $\lr {S,S}$ and therefore 
we have that  
$S^{-1} \mathbf 1 \eqv B \lr {S,S}$
is contractible.
\end{remark}

\begin{definition}[invertible action]
We say that $S$ acts \emph{invertibly} if each translation $s \boxtimes : X \to X$ is a homotopy equivalence.
\end{definition}

If $S$ acts invertibly on $S^{-1}X$ by $s \boxtimes (t,x) = (s \square t, x)$, which the homotopy inverse $s \boxtimes (t,x) = (t,s \boxtimes x)$
since there is a natural transformation $(t,x) \mapsto (s \square t, s \boxtimes x)$.

We combine these obesrvations to get the following,
\begin{proposition}
Since $\pi_0S$ is a multiplicatively closed subset of the ring $H_0(S) = \Z[\pi_0S],$
it acts on $H_*(X)$ and acts invertibly on $H_*(S^{-1}X)$.
The functor $X \to S^{-1}X$ described by $x \mapsto (0,x)$ induces the map

\begin{equation}
  (\pi_0S)^{-1} H_q(X) \to H_q(S^{-1}X)
  \label{eq:localaction}
\end{equation}
\end{proposition}

We have the following result due to Quillen, showing us that $B(S^{-1}S)$ is indeed a group completion of $BS$. (See [Grayson])
\begin{theorem}[Quillen]\label{thm:BSgroup}
Let $S$ be a faithful symmetric monoidal groupoid. Then (\ref{eq:localaction}) is an isomorphism for all $X$ and $q$.
In particular, $B(S^{-1}S)$ is a group completion of the $H$-space $BS$.
\end{theorem}
\subsection{Relation to the $+$-construction}

In order to deduce $K(\mathbf P(R))$, we will first deduce $K(\mathbf F(R))$ and use a fact about cofinality to 
get our result.

\begin{theorem}[Quillen]
Let $S = \coprod GL_n(R)$ which is equivalent to $\mathbf F(R)$.
Then $B(S^{-1}S)$ is the group-completion of $BS$ and 
\[
B(S^{-1}S) \eqv \Z \times BGL(R)^+
\]
\end{theorem}
\begin{proof}
We will try to find an acyclic map from $BGL(R)$ to the connected basepoint component of $B(S^{-1}S)$.
We perform this via the mapping telescope construction which describes the homotopy colimit.
\[
BGL(R) = \varhocolim BGL_n(R)
\]
Homotopy classes of maps $BGL(R) \to B(S^{-1}S)$ correspond to compatible homotopy classes maps
$BGL_n(R) \to B(S^{-1}S)$ for all $n$.
Note that group maps $GL_n(R) \to \Aut_{S^{-1}S}(R^n, R^n)$ describe maps from $BGL_n(R) \to B(S^{-1}S)$.
In fact, we have maps $f_n(g) = (g,1)$ for each $n$ and the following diagram commutes.
\[
\begin{tikzcd}
GL_n(R) \arrow[r, "f_n"] \arrow[d, "\square R"'] & \Aut(R^n,R^n) \arrow[d, "\square (R R)"] \\
GL_{n+1}(R) \arrow[r, "f_{n+1}"] & \Aut(R^{n+1}, R^{n+1})
\end{tikzcd}
\]
This implies we have natural transformations $f_{n} \To f_{n+1}(\square R)$,
which induce the homotopies $f_{n} \To f_{n+1}(\square R) : BGL_n(R) \to B(S^{-1}S)$.
Due to this compatibility of $f_n$, we get a map out of the mapping telescope
$f : BGL(R) \to B(S^{-1}S)$.

Furthermore, the map lands in the connected component $Y_S$ of the identity in $B(S^{-1}S)$,
which is also a $H$-space. 

We will now show that this map $BGL(R) \to Y_S$ is an isomorphism on integral homology. 
By Proposition \ref{prop:acyclic_H_space}, this will induce a homotopy equivalence $BGL(R)^+ \to Y_S$.

Let $e \in \pi_0BS$ denote the path connected component of $R$ in $BS$.
By Theorem \ref{thm:BSgroup}, $H_*(BS^{-1}S)$ is the localization of the ring $H_*(BS)$ at $\pi_0S = \{e^n\}$.
But this is just the colimit
\[
H_*(BS) \to H_*(BS) \to H_*(BS) \to H_*(BS) \to \cdots 
\]
where the maps are given by the translation $\oplus R : S \to S$.
So $H_*B(S^{-1}S) \cong H_* (Y_S) \otimes \Z[e,e^{-1}]$ which will allow us to recover $B(S^{-1}S) \eqv \Z \times Y_S$.
Finally, we have the chain
\[H_*(Y_S) \cong \colim H_*(BGL_n(R)) = H_*(BGL(R))\]
And this shows $H_*(BGL(R)) \to Y_S$ is a homology isomorphism.
\end{proof}

\subsection{Cofinality and Equivalence with $+$-construction}

A monoidal functor $f : S \to T$ of faithful symmetric monoidal groupoids is said to be \emph{cofinal} if for 
every $t$ in $T$, there is some $s$ in $S$ and $u$ in $T$ such that $t \square u \cong f(s)$.
The inclusion functor $\mathbf F(R) \to \mathbf P(R)$ for instance, is cofinal,
because every projective module is a direct summand of a free module.

We have the following result which allows us to $K_n(T)$, a proof of which is outlined in [Weibel].

\begin{theorem}[Cofinality Theorem]
Let $f : S \to T$ be a cofinal monoidal functor of symmetric monoidal categories. Then
\begin{enumerate}
  \item If $T$ acts on $X$, then $S$ acts on $T$ via $f$, and $S^{-1}X \eqv S^{-1}X$.
  \item If $\Aut_S(s) \cong \Aut_T(fs)$ for all $s$ in $S$, then the basepoint components
  of $K(S)$ and $K(T)$ are homotopy equivalent.
  In particular, $K_n(S) \cong K_n(T)$ for $n \geq 1$.
\end{enumerate}
\end{theorem}
%\begin{proof}
%By cofinality, $S$ acts invertibly on $X$ if and only if $T$ acts invertibly on $X$. 
%
%\end{proof}

Owing to this Theorem, the fact that $\mathbf F(R)$ is cofinal in $\mathbf P(R)$, 
we have the following corollary.

\begin{corollary}
Let $S = \iso \mathbf P(R)$ be the (faithful) symmetric monoidal groupoid of finitely generated
projective $R$-modules. Then
\begin{equation}
B(S^{-1}S) \eqv K_0(R) \times BGL(R)^+
  \label{eq:Sequals+}
\end{equation}
\end{corollary}

%\begin{theorem}
%%Let $S$ be a symmetric monoidal groupoid with faithful transitions satisfying cofinality.
%
%We have then 
%\[
%K^\square(S) \eqv K^\square_0 (S) \times B \Aut(S)^+A 
%\]
%\end{theorem}

\section{Quillen's $Q$-construction}\label{sec:Qconstruction}
In [Quillen, 19XX], Quillen proposed another construction for the $K$-theory space.
Given an exact category $\A$, []
we will define a category $Q \A$, and then define $K_n(\A) = \pi_{n+1}(BQ\A)$ for all $n \geq 0$.

\begin{definition}[$Q$-construction]
Let $\A$ be an exact category. We define $Q\A$ as the category with the same objects as $\A$, and the morphisms between $C \to C'$ are 
equivalence classes of diagrams
\[
\begin{tikzcd}
C & D \arrow[l,two heads, "j"] \arrow[r,tail, "i"] & C' 
\end{tikzcd}
\]
where two such diagrams are considered isomorphic when there is an isomorphism of these diagrams which is identity on $C$ and $C'$.

The composition of the above arrow with $\Qmor{C'}{D'}{C''}$ is given by $\Qmor{C}{D \times_{C'} D'}{C''}$ from the diagram
\[
\begin{tikzcd}%[column sep=large]
D \times_{C'} D' \arrow[r, tail, "\pi_2"] \arrow[d, two heads, "\pi_1"'] & D' \arrow[r, tail, "i'"] \arrow[d, two heads, "j'"]& C''\\
D \arrow[r, tail, "i"'] \arrow[d, two heads, "j"'] & C'\\
C
\end{tikzcd}
\]
It is clear that composition is well-defined and associative and so $Q\A$ forms a well-defined category.
\end{definition}

[Set theoretic remark]

\begin{remark}
Admissible monomorphisms $C \admono C'$ determine a morphism $C \to C'$ by taking $D = C'$ and $D \adepi C'$ as identity.
We will call these morphisms \emph{injective}.

Similarly, admissible epimorphisms $C \adepi C'$ determine a morphism $C \to C'$ by taking $D = C$ and $C \admono D$ as identity.
We will call these morphisms \emph{surjective}.

Note that these distinguished families of morphisms are closed under composition, and generate all morphisms.
Furthermore, any morphism $\Qmor{C}{D}{C'}$ factors as $i_*j^*$ where 
$i_*$ is the injective morphism determined by $D \admono C'$ and $j^*$ is the surjective morphism determined by $C \adepi D$.
This factorization is unique up to unique isomorphism.
\end{remark}

\subsection{Alternative Description of Morphisms in $Q\A$}
There is an alternative way to describe the morphisms $C \to D$ in $Q\A$
in terms of \emph{subobjects} and \emph{subquotients}. Fix some object $C'$ in $\A$.

\begin{definition}
An \emph{admissible subobject} of $C'$ is an equivalence class of admissible monomorphisms
$C \admono C'$ with equivalence given by whether they factor through each other, 
or in other words as isomorphism in $\A \back C$. 

An \emph{admissible subquotient} is defined in the dual manner, using admissible epimorphisms.
\end{definition}

One can define a partial order on the set of admissible monomorphisms into $C'$.
We declare $C_1 \leq C_2$ if there is an admissible monomorphism $C_1 \admono C_2$, which if it exists
has to be unique.
We then say $(C_1, C_2)$ is an \emph{admissible layer} in $C'$ and the cokernel $C_2/C_1$ is an \emph{admissible subquotient} of $C'$.

\begin{proposition}\label{prop:Qadpair}
The morphisms from $C \to C'$ may be identified with pairs $((C_1,C_2),\varphi)$,
where $(C_1,C_2)$ is an admissible layer in $C'$ and $\varphi : C \xrightarrow{\cong} C_2/C_1$ is an isomorphism in $\A$.
\end{proposition}
\begin{proof}
For $\Qmor{C}{D}{C'}$, take the admissible layer $(C_1, C_2)$ where $C_2 = D$ and $C_1$ is the kernel of $C \adepi D$.
Clearly, since $C_1$ gives a subobject of $D = C_2$ which gives a subobject of $C'$, the given pair is indeed 
an admissible layer. A unique isomorphism follows since both $C$ and the $C_2/C_1$ are cokernels of the kernel of an admissible epimorphism.

\[
\begin{tikzcd}
  & \ker(j) \arrow[d, tail] \arrow[dl, "0"'] \arrow[dr, tail] & \\
C\arrow[dr, "\exists ! \cong"'] & D \arrow[l, two heads, "j"] \arrow[r, tail, "i"'] \arrow[d, two heads] & C' \\
  & D/\ker(j)&
\end{tikzcd}
\]
Conversely, an admissible layer $(C_1,C_2)$ in $C'$ evidently determines a morphism 
$\Qmor{C_2/C_1}{C_2}{C}$ in $Q\A$ where the isomorphism may be taken as identity.

Furthermore, composition of morphisms $C \to C'$ and $C' \to C''$ may be described this way and arises from the fact that
subquotient of a subquotient is a subquotient. 

Let morphisms $C \to C'$ and $C' \to C''$ correspond to $((C_1,C_2),\varphi)$ and $((C_1',C_2'),\psi)$
where $(C_1,C_2)$ 
and $(C_1',C_2')$ are admissible layers in $C$ and $C''$ respectively with isomorphisms
$\varphi : C \xrightarrow{\cong} C_2/C_1$ and $\psi : C' \xrightarrow{\cong} C_2'/C_1'$.
Let the map $C_2' \to C_2'/C_1'$ be denoted $q$.

Then the composite $C \to C' \to C''$ corresponds to the layer $(D_1,D_2)$
where $D_1 = C_1 \times_{C'} C_2'$ and $D_2 = C_2 \times_{C'} C_2'$
with the isomorphism $\Phi : C \xrightarrow{\cong} D_2/D_1$ given by
\[
  C \cong C_2/C_1 
\]
This follows from the evident isomorphisms and the fact that 
$D_2$ is the pullback of $C_2 \admono C' \adepi C_2'$.
\end{proof}
This characterization is particularly helpful to eludicate certain properties of the category $Q\A$.

\begin{example}
Using this characterization, we can describe the morphisms with source $0$ in $Q\A$.
Morphisms $0 \to C$ are in 1-1 correspondence with the subobjects of $C$, thus injective. 
There is a canonoical morphism $C \to 0$ since $0$ has a unique subquotient. 
\end{example}
\begin{example}
Isomorphisms $C \xlongrightarrow{\cong} D$ in $Q\A$ are in 1-1 correspondence with isomorphisms $C \xlongrightarrow{\cong} D$ in $\A$.
\end{example}

\subsection{$K$-Theory using the $Q$-construction}

We now prove a few homotopy-theoretic properties of the $Q$-construction,
before we eventually define $K_n$-groups corresponding to a given (skeletally small) exact category.

\begin{proposition}
An exact functor $F : \A \to \B$ induces a functor $QF : Q\A \to Q\B$.
\end{proposition}
\begin{proof}
On the level of objects, it suffices to define $QF$ precisely as $F$.
Suppose $g : C \to D$ is a morphism in $Q\A$. We may identify $g$ with an admissible layer $(D_1,D_2)$ of $D$. 
Define $QF(g)$ to be the morphism $FC \to FD$ in $Q\B$ corresponding to the admissible layer $(FD_1,FD_2)$ of $FD$,
which is indeed an admissible layer since the exactness relations in $\A$ carry over to $\B$
by the exact functor $F$.
Functoriality follows since compositions are defined naturally across $Q\A$ and $Q\B$.
\end{proof}

We present a different proof of the theorem proven by Quillen in [].

\begin{theorem}[Quillen]
Let $\A$ be a small exact category.
The fundmental group of $BQ\A$ at the basepoint $0$ is canonically isomorphic to 
the $K_0(\A)$ group defined in Chapter 2.
\end{theorem}
\begin{proof}
Note that $\pi_1(BQ\A)$ is generated by loops given by the morphisms $0 \to C \back 0$
for each $C$ in $Q\A$, where $0 \to C$ is the injective morphism given by the subobject $0 \admono C$,
and $C \back 0$ is the surjective morphism given by $0 \adepi C$.
This also shows $BQ\A$ is (path) connected.

It is clear that the loops given by $0 \to C \back 0$ determine $C$ up to isomorphism,
and loops $0 \to C \back 0$ generate $\pi_1(BQ\A)$.

Given any short exact sequence $C' \admono C \adepiback C''$ in $\A$, we wish to show that the loop given by
the concatenated diagram $0 \to C' \back 0 \to C'' \back 0$ in $Q\A$ is homotopic to the loop given by $0 \to C \back 0$.
Drawing out the morphisms of these diagrams in $Q\A$, we have the following diagram in $\A$
\[
\begin{tikzcd}
0 & C' \arrow[l, two heads] \arrow[r, tail] & C' & 0 \arrow[r, two heads] \arrow[l, tail] & 0 \\
  & 0  \arrow[ul,two heads] \arrow[dr,tail] & C'\arrow[u,two heads] \arrow[d,tail] & 0  \arrow[dl, tail] \arrow[ur,two heads]  & 0 \arrow[u, two heads] \arrow[d, tail]\\
  &             & C          & C \arrow[l, tail] \arrow[r, two heads] & C'' \\
  &             &            & C \arrow[ul, tail] \arrow[dr,two heads] & C''  \arrow[u, tail] \arrow[d, two heads]\\
  &             &            &    & 0
\end{tikzcd}
\]
Note that there are four triangles with vertices and morphisms in $Q\A$ which divide this diagram as displayed here. 
To show the desired claim, it suffices to show that each of the four traingles commute in $Q\A$,
which would imply each corresponds to a $2$-simplex in $\pi_1(BQ\A)$, making the loop
$0 \to C' \back 0 \to C'' \back 0$ homotopic to $0 \to C \back 0$ due to shared edges.

But in each of the triangles the computation of the pullback in the composition is trivial 
and yields either the injective morphism $0 \to C$ (by $\Qmor 0 0 C$) or the surjective morphism $0 \to C$ (by $\Qmor 0 C C$).
In any case, it is equal to the composite edge.

It also follows that $\pi_1(BQ\A)$ is abelian.

We can now define a group homomorphism $\Theta : K_0(\A) \to \pi_1(BQ\A)$ by
mapping $[C]$ to the loop $\Qloopz C$ as mentioned. This is a (surjective) group homomorphism
since exact sequences lead to homotopic loops. 

To show $\Theta$ is injective, 
We construct an $\Phi$ from morphisms in $Q\A$ to $K_0(\A)$ by mapping
the morphism determined by the layer $(D_1,D_2)$ of $C'$ to $D_1$.
This is additive since 
\end{proof}

[Set Theory remark]

\begin{definition}
Let $\A$ be a small exact category. Then $K\A$ is defined to be the space $\Omega BQ\A$ pointed at $0$
and further
\[
K_n(\A) := \pi_n(K\A) = \pi_{n+1}(BQ\A)
\]
\end{definition}

This definition agrees with $K_0(\A)$ defined in Chapter 2. Further, any exact functor $F : \A \to \B$
induces a functor $Q\A \to Q\B$, hence $BQ\A \to BQ\B$ showing that $K_n(\A) \to K_n(\B)$
showing that the $Q$-construction leads to a functorial $K$-theory.

\begin{definition}
If $R$ is a unital ring,
we define $K(R) := K(\mathbf P(R))$ and consequently
\[K_n(R) = K_n(\mathbf P(R)).\] 
This matches our original definition for $n = 0$.
We prove in \S\ref{sec:+=Q} the $+=Q$ Theorem, which implies that this matches with the description of the 
$K_n$ groups defined using the $+$-construction in \S\ref{sec:BGL(R)^+}.
\end{definition}

\begin{definition}
If $R$ is a Noetherian ring with unity,
we define $G(R) := K(\mathbf M(R))$ and consequently 
\[G_n(R) := K_n(\mathbf M(R)).\] 
This gives us $G$-Theory of unital Noetherian rings and matches the classical definition for $n = 0$.
\end{definition}


%\subsection{Functoriality of the $Q$-construction}

\subsection{Key Properties of the $Q$-construction}

\subsubsection{Morita Invariance}
Since $K$-theory from the $Q$-construction is functorial,
for Morita-equivalent rings $R, S$, i.e., there is an equivalence $\mathbf P(R) \eqv \mathbf P(S)$
of categories, from which it follows that $K_n(R) \cong K_n(S)$.

\subsubsection{Cofinality}

\section{The $``+=Q"$ Theorem} \label{sec:+=Q}

Given an additive category $\A$, there are two approaches to defining $K$-theory over $\A$.
One is through

\begin{theorem}[Quillen]\label{thm:Quillen}
Let $\mathcal A$ be a split-exact category 
and $S = \iso \mathcal A$ be the symmetric monoidal groupoid of isomorphisms in $\A$ with $\square = \oplus$,
then \[\Omega BQ \mathcal A \eqv B(S^{-1}S) \]
\end{theorem}

\begin{corollary}[$+=Q$]\label{thm:+=Q}
For any unital ring $R$,
\[\Omega B Q \mathbf P (R) \eqv K_0(R) \times BGL(R)^+\]
This means that the $K$-Theory of $\mathbf P(R)$ from the $Q$-construction is equivalent to the 
$K$-Theory defined using the $+$-construction.
\end{corollary}

\subsection{Proof of the $+=Q$ Theorem}

For the proof,
we wish to produce a fibration sequence
\[
S^{-1}S \to \textbf{??} \to Q\A
\]
where \textbf{??} is chosen appropriately 
and $B(\textbf{??})$ is contractible.
Then upon taking the homotopy long exact sequence of the fibration,
we would have our desired homotopy equivalence $\Omega BQ\A \eqv BS^{-1}S$.
This sequence of leading up to the theorem is taken from [Weibel].

\begin{definition}%[(Give it a name?)]
Given an exact category $\A$, we define $\E\A$ where 
the objects of $\E\A$ are admissible exact sequences in $\A$,
and a morphism between sequences $\adexact{C'}{C}{C''}$ and $\adexact{D'}{D}{D''}$
is an equivalence class of the following diagrams admssible exact sequences $\adexact{D'}{C}{E}$
such that the following diagram commutes
\[
\begin{tikzcd}
C' \arrow[r,tail] & C \arrow[r,two heads] \arrow[d,equal] & C''\\
D' \arrow[r,tail] \arrow[u,tail, "\alpha"]  \arrow[d,equal] & C \arrow[d,tail, "\beta"'] \arrow[r,two heads] & E\arrow[d,tail] \arrow[u,two heads] \\
D' \arrow[r,tail] & D \arrow[r,two heads] & D''
\end{tikzcd}
\]
Two such diagrams are considered equivalent when there is an isomorphism between them which is
identity at all vertices except for $E$.
\end{definition}

%\begin{remark}
%In the language of subobjects and subquotients, $\E\A$ consists of an object $C''$
%along with a jk
%\end{remark}

Note since the right edge in the diagram describes a morphism in $Q\A$,
we have a functor $f : \E\A \to Q\A$ given by picking $C''$ from 
the exact sequence $\adexact{C'}{C}{C''}$.

We now consider the fiber category $f^{-1}(E)$ for an object $E$ which we denote $\E_E$.

\begin{remark}
If we consider the identity morphism to be fixed at $E$ , then the fiber category
$\E_E$ consists of exact sequences with target $E$, and the morphisms from 
$\adexact{C'}{C}{E}$ to $\adexact{D'}{D}{E}$ 
are diagrams of the form
\[
\begin{tikzcd}
C' \arrow[r,tail] & C \arrow[r,two heads] \arrow[d, "\cong", "\beta"'] &E \arrow[d,equal] \\ 
D' \arrow[r,tail] \arrow[u, "\cong"', "\alpha"] & D \arrow[r,two heads] &E 
\end{tikzcd}
\]
It follows automatically that each morphism in $\E_E$ is an isomorphism which we may describe as a pair
$(\alpha, \beta)$ where $\alpha : D' \cong C'$ and $\beta : C \cong D$.
\end{remark}

The following result allows us to relate the fiber category with the symmetric monoidal groupoid $S = \iso \A$.

\begin{proposition}
The fiber category $\E_0$ is homotopy equivalent to $S = \iso \A$.
\end{proposition}
\begin{proof}
Let $F : S \to \E_0$ be the functor given by mapping an object $C$ to $\adexact{C}{C}{0}$,
and an isomorphism $f : C \to C$ to the morphism $(f, f^{-1})$ with the diagram 
\[
\begin{tikzcd}
C \arrow[r,tail] & C \arrow[r,two heads] \arrow[d, "f^{-1}"'] &E \arrow[d,equal] \\ 
C \arrow[r,tail] \arrow[u, "f"] & C \arrow[r,two heads] &E 
\end{tikzcd}
\]
In fact, since every object in $\E_0$ is naturally isomorphic to such a sequence,
the functor $F$ is defined naturally. It is also a full embedding. This gives us the claim:
\end{proof}

\begin{proposition}
For any object $E$ in $\A$, $\E_E$ is a symmetric monoidal category,
and there is a (faithful) monoidal functor $\eta_E : S \to \E_E$ sending $C$ to the
sequence $\adexact{C}{C\oplus E}{E}$.
\end{proposition}
\begin{proof}
We define the monoidal structure on $\E_E$ as follows. 
Let $e = (\adexact{0}{E}{E})$ be the identity.
Let $x$ denote $\adexact{C'}{C}{E}$ and $y$ denote $\adexact{D'}{D}{E}$.
Then we define $x \square y$ as
\[
\adexact{C' \oplus D'}{C \times_E D}{E}
\]
It follows that this obeys the necessary unit, commutatitivity and coherence laws
and does define a symmetric monoidal structure on $\E_E$.
The symmetric monoidal functor $S \to \E_E$ follows naturally by mapping
$C$ to the split-exact sequence $\adexact{C}{C \oplus E}{E}$.
[]
It is also faithful. 
\end{proof}

Now, here is a wonderful observation which reduces the complexity of the fiber category $\E_E$ in a restricted case. 
\begin{remark}
If $\A$ were split-exact, then every sequence in $\E_E$ is of the form $\adexact{C}{C \oplus E}{E}$,
hence comes from $S$ via the functor $\eta_E$. 
Consequently, the category $\langle S, \E_E \rangle$ is connected.
\end{remark}

\begin{proposition}
If $\A$ is split exact, each $S^{-1}S \to S^{-1}\E_E$ is a homotopy equivalence.
\end{proposition}
\begin{proof}
Note the sequence $S^{-1}S \to S^{-1}\E_E \to \langle S, \E_E \rangle$ is a fibration.
It suffices to show that $L = \langle S, \E_E \rangle$ is contractible.
Note that the monoidal operation on $\E_E$ induces a monoidal operation on $L$,
making $BL$ a $H$-space. Since $L$ is connected, by [Whitehead] it implies $BL$ is group-like
and has a homotopy inverse.

Now for any exact sequence $x = (\adexact{C}{C \oplus E}{E})$ we have a diagonal morphism 
$\Delta_x : x \to x \square x$ in $L$,
\[
\begin{tikzcd}
C     \arrow[r,tail] \arrow[d] & D       \arrow[r,two heads] \arrow[d] & C \arrow[d, equal] \\
C \oplus C \arrow[r,tail]           & D \times_C D \arrow[r,two heads]           & C 
\end{tikzcd}
\]
The diagonal $\Delta_x$ is natural in $x$. 
This natural transformation induces a homotopy between $\id : BL \to BL$ and $(-)^2 : BL \to BL$. 
Using the homotopy inverse to subtract away $\id$, we get $\id : BL \to BL$ is null-homotopic. 

Hence $BL$ is contractible.
\end{proof}

\begin{lemma}\label{lemma:Qbasechange}
For any morphism $f : E' \to E$ in $Q\A$, there is a canonical ``base-change" functor
$f^* : \E_E \to \E_{E'}$ and a natural transformation $\eta : f^* \Rightarrow i$,
where $i$ is the inclusion $\E_E \hookrightarrow \E\A$.
\end{lemma}

\begin{theorem}\label{thm:splitfibration}
If $\A$ is a split exact category and $S = \iso \A$, then the sequence
\[S^{-1}S \to S^{-1}\E\A \to Q\A\]
is a homotopy fibration.
\end{theorem}
\begin{proof}
We use Quillen's Theorem B to prove this. 
For Theorem B to apply, we need to show that the base changes
$f^*$ are homotopy equivalences.

We show the same on properly injective and properly surjective morphisms $0 \to C$.
For properly injective morphisms, the base change is given by $(\adexact{C}{D}{E}) \mapsto (\adexact{C}{C}{0})$
which can be identified with $C$ in $S$.
The composition $S^{-1}S \to S^{-1}\E_E \to Q\A$ is thus identity.

For properly surjective morphisms, the base change is $(\adexact{C}{D}{E}) \mapsto (\adexact{D}{D}{-1})$ which
can be identified with $D$ in $S$.
The composition $S^{-1}S \to S^{-1}\E_E \to Q\A$ in this case is $C \mapsto (D = C \oplus E)$,
which is realized by the natural transformation given by $E \oplus $,
hence is also a homotopy equivalence.

Hence we can conclude that in general, $f^*$ is a homotopy equivalence.
\end{proof}

%\begin{proposition}
%The category $S^{-1}\E\A$ is contractible.
%\end{proposition}

\begin{proof}[\textbf{Proof of Theorem \ref{thm:Quillen}}]
The proof follows from Theorem \ref{thm:splitfibration} if we show that $S^{-1}\E\A$ is contractible.
Note that $\E\A$ is contractible. [] Any action of $S$ on a contractible category must be invertible. [] %\ref
So we have the homotopy equivalence $\E\A \to S^{-1}\E\A$, concluding that $S^{-1}\E\A$ is contractible.
\end{proof}

This concludes the section and gives us Theorem \ref{thm:Quillen} and Corollary \ref{thm:+=Q}. 
